% SHARD-ID: AQM-08-SYMPLECTIC-CSS-BRIDGE
% SHARD-TITLE: Chain Complexes as Symplectic CSS Stabilizers
% SHARD-SUMMARY: Derives the exact relation between chain/cochain CSS data and isotropic stabilizer subspaces.
% SHARD-SUMMARY: Identifies the logical Pauli module as the symplectic reduction \(L^\perp/L\).
% SHARD-SUMMARY: Records the prime-field qudit validation run for the oriented toric complex.
% SHARD-KEYWORDS: QECC, symplectic geometry, CSS codes, chain complex, qudits, stabilizer

\section{Chain Complexes as Symplectic CSS Stabilizers}
\label{sec:symplectic-css-bridge}

This section answers the symplectic-geometry question. The conclusion is:

\begin{quote}
The chain/cochain CSS construction is not an alternative to the symplectic
description. It is a functorial way to manufacture an isotropic stabilizer
subspace \(L\) inside the Pauli symplectic vector space. The usual homology and
cohomology groups appear as the two halves of the symplectic reduction
\(L^\perp/L\).
\end{quote}

The local source anchors are as follows. Gottesman's TeX source records CSS
commutation as row orthogonality in
\path{references/symplectic_qecc/gottesman_9705052_source/Thesis.tex},
lines 1227--1246, and the binary symplectic stabilizer form in lines
1293--1341. Ashikhmin--Knill's local TeX source
\path{references/symplectic_qecc/ashikhmin_knill_0005008_source/n9.tex},
lines 249--313, gives the nonbinary shift/phase commutator; lines 351--407
relate nonbinary stabilizers to self-orthogonal classical data. The
Bombin--Martin-Delgado local source
\path{references/toric_code/bombin_martin_delgado_0605094_source/HomologicalErrorCorrection6.tex},
lines 1206--1259, writes qudit Pauli operators with a symplectic commutator;
lines 2050--2136 define the chain/cochain pairing; lines 2383--2452 state the
homological-code construction as a symplectic code.

\subsection{The Symplectic Stabilizer Side}

First take prime qudit dimension \(p\). Put one qudit on each of \(n\) sites.
The Pauli label space is
\[
  V=\mathbb F_p^n\oplus\mathbb F_p^n .
\]
We write an element as \((x,z)\), where \(x\) is the vector of \(X\)-exponents
and \(z\) is the vector of \(Z\)-exponents. With
\(\zeta=\exp(2\pi i/p)\), define
\[
  X^x|a\rangle=|a+x\rangle,\qquad
  Z^z|a\rangle=\zeta^{z\cdot a}|a\rangle ,
\]
where \(a,x,z\in\mathbb F_p^n\). Then
\[
  X^xZ^zX^{x'}Z^{z'}
  =
  \zeta^{z\cdot x'-x\cdot z'}X^{x'}Z^{z'}X^xZ^z .
\]
Thus the commutator is controlled by the symplectic form
\[
  \omega((x,z),(x',z'))=z\cdot x'-x\cdot z' .
\]
This is the local convention used throughout the section: the first component
is the \(X\)-exponent and the ordered \(X\)-then-\(Z\) mixed pairing is negative.

A stabilizer specified by exponent labels \(L\subset V\) is commuting exactly
when
\[
  \omega(\ell,\ell')=0\quad\text{for all }\ell,\ell'\in L.
\]
That is, \(L\) is isotropic. Exponent labels determine this centralizer and
commutation condition; an actual stabilizer code also requires compatible
phases/eigenvalue characters. For the canonical CSS lifts used below, if
\(\dim L=r\), the joint eigenspace has dimension
\[
  p^{n-r}.
\]
The logical Pauli labels are not \(V/L\), because Pauli operators must preserve
the code space. They are
\[
  L^\perp/L,\qquad
  L^\perp=\{v\in V:\omega(v,\ell)=0\text{ for all }\ell\in L\}.
\]
This is the finite symplectic reduction. Its dimension is \(2(n-r)\).

\subsection{The CSS Case Inside the Symplectic Module}

A CSS stabilizer is specified by two subspaces
\[
  R_X,R_Z\subset \mathbb F_p^n.
\]
The \(X\)-type labels are \((x,0)\) with \(x\in R_X\), and the \(Z\)-type labels
are \((0,z)\) with \(z\in R_Z\). Hence
\[
  L=(R_X\oplus 0)+(0\oplus R_Z)\subset V.
\]
The same-type commutators vanish automatically. The mixed commutator is
\[
  \omega((x,0),(0,z))=-x\cdot z .
\]
Therefore
\[
  L\text{ isotropic}
  \quad\Longleftrightarrow\quad
  x\cdot z=0\ \text{for all }x\in R_X,\ z\in R_Z.
\]
If \(H_X\) and \(H_Z\) are matrices whose rows span \(R_X\) and \(R_Z\), this
condition is exactly
\[
  H_XH_Z^T=0.
\]
So the CSS row-orthogonality condition is literally the isotropy condition in
the Pauli symplectic vector space.

\subsection{Feed in a Chain Complex}

Now let
\[
  C_2\xrightarrow{\partial_2}C_1\xrightarrow{\partial_1}C_0
\]
be a finite chain complex over \(\mathbb F_p\), so
\[
  \partial_1\partial_2=0.
\]
Put the physical qudits on the chosen basis of \(C_1\), so \(n=\dim C_1\). Let
\[
  \delta^0=\partial_1^T:C^0\to C^1,\qquad
  \delta^1=\partial_2^T:C^1\to C^2 .
\]
After identifying \(C^1\) with coordinate row vectors on the edge basis, define
\[
  R_X=\operatorname{im}\delta^0,\qquad
  R_Z=\operatorname{im}\partial_2.
\]
Thus the stabilizer subspace is
\[
  L=(\operatorname{im}\delta^0\oplus 0)
    +(0\oplus \operatorname{im}\partial_2)
  \subset C^1\oplus C_1 .
\]

We now derive isotropy, not just state it. Take \(\alpha\in C^0\) and
\(\beta\in C_2\). The corresponding \(X\)- and \(Z\)-type labels are
\((\delta^0\alpha,0)\) and \((0,\partial_2\beta)\). Their symplectic pairing is
\[
\begin{aligned}
  \omega((\delta^0\alpha,0),(0,\partial_2\beta))
  &=-\langle \delta^0\alpha,\partial_2\beta\rangle       \\
  &=-\langle \alpha,\partial_1\partial_2\beta\rangle      \\
  &=0 .
\end{aligned}
\]
The first line is the Pauli symplectic form. The second line is the definition
of the coboundary as the transpose of the boundary. The last line is exactly
\(\partial_1\partial_2=0\). This is the clean bridge: the chain-complex
identity is the CSS isotropy identity.

\subsection{The Exact Sequence Hidden in Symplectic Geometry}

The stronger statement identifies the symplectic reduction. Let
\[
  B^1=\operatorname{im}\delta^0,\quad
  Z^1=\ker\delta^1,\quad
  B_1=\operatorname{im}\partial_2,\quad
  Z_1=\ker\partial_1.
\]
Then \(L=B^1\oplus B_1\), using the CSS embedding into \(C^1\oplus C_1\).

Compute \(L^\perp\). A general Pauli label is \((x,z)\in C^1\oplus C_1\). It
is orthogonal to every \((u,0)\) with \(u\in B^1\) iff
\[
  0=\omega((x,z),(u,0))=z\cdot u
  \quad\text{for all }u\in \operatorname{im}\delta^0.
\]
This says \(z\) annihilates \(\operatorname{im}\delta^0\), equivalently
\[
  z\in(\operatorname{im}\delta^0)^\perp
  =\ker\partial_1=Z_1 .
\]
It is orthogonal to every \((0,b)\) with \(b\in B_1\) iff
\[
  0=\omega((x,z),(0,b))=-x\cdot b
  \quad\text{for all }b\in \operatorname{im}\partial_2.
\]
This says
\[
  x\in(\operatorname{im}\partial_2)^\perp
  =\ker\delta^1=Z^1 .
\]
Therefore
\[
  L^\perp=Z^1\oplus Z_1 .
\]
Quotienting by \(L=B^1\oplus B_1\) gives the short exact sequence
\[
  0\longrightarrow B^1\oplus B_1
  \longrightarrow Z^1\oplus Z_1
  \longrightarrow H^1\oplus H_1
  \longrightarrow 0,
\]
and hence
\[
  L^\perp/L\cong H^1\oplus H_1.
\]
This is the precise connection. The chain/cochain short exact sequence is the
CSS presentation of the symplectic reduction of an isotropic stabilizer
subspace.

\subsection{Lagrangians and Logical Bases}

If \(L\) has dimension \(n-k\), then \(L^\perp/L\) has dimension \(2k\). The
code encodes \(k\) qudits. A stabilizer code itself is therefore represented by
an isotropic subspace \(L\), not usually a Lagrangian subspace of \(V\). It
becomes Lagrangian only when \(k=0\), i.e. for a stabilizer state.

The Lagrangians appear one level down, in the logical symplectic module
\[
  L^\perp/L\cong H^1\oplus H_1.
\]
For example,
\[
  \Lambda_X=Z^1\oplus B_1,\qquad
  \Lambda_Z=B^1\oplus Z_1
\]
are isotropic subspaces of \(L^\perp\) containing \(L\). Their quotients are
\[
  \Lambda_X/L\cong H^1,\qquad
  \Lambda_Z/L\cong H_1.
\]
These are complementary Lagrangian subspaces of the logical Pauli module, and
the induced symplectic pairing on the ordered \(X\)-then-\(Z\) logical labels
is the negative of the natural evaluation pairing:
\[
  H^1\times H_1\longrightarrow \mathbb F_p,\qquad
  ([x],[z])\mapsto -x(z).
\]
Equivalently, if the arguments are ordered \(Z\)-then-\(X\), the same
symplectic form gives the positive evaluation \(x(z)\).
Choosing a logical computational basis is choosing one of these Lagrangian
polarizations.

\subsection{Qudit Generalization}

For prime \(p\), the derivation above is literal. For prime powers \(q=p^m\),
replace \(\mathbb F_p\) by \(\mathbb F_q\) and use the additive character
\[
  \chi(t)=\exp(2\pi i\,\operatorname{Tr}_{\mathbb F_q/\mathbb F_p}(t)/p).
\]
The Pauli commutator is then controlled by the trace-symplectic form
\[
  \operatorname{Tr}_{\mathbb F_q/\mathbb F_p}
  \bigl(z\cdot x'-x\cdot z'\bigr).
\]
The CSS chain-complex condition over \(\mathbb F_q\) is still
\[
  \partial_1\partial_2=0,
\]
and the same calculation proves isotropy. Thus the chain/cochain construction
does generalize cleanly to field-valued qudits.

The caveat is composite non-field dimension \(D\). One can work over
\(\mathbb Z_D\), but then the clean vector-space statements about rank,
dimension, and duality become module statements. Torsion can matter. This is
why the field symplectic formulation is the robust one: the chain complex is
safe when interpreted as producing the isotropic subspace \(L\), while the
logical information is recovered from the symplectic reduction \(L^\perp/L\).

\subsection{Exact Validation}

The validation script is
\path{scripts/lattice_codes/symplectic_css_bridge_validation.jl}. It builds
oriented boundary matrices for the \(k=4\) periodic square torus over
\(\mathbb F_p\) for \(p=2,3,5\). It
verifies:
\[
  \partial_1\partial_2=0,\qquad L\subseteq L^\perp,
\]
\[
  \operatorname{rank}\partial_1=15,\qquad
  \operatorname{rank}\partial_2=15,\qquad
  \dim H_1=32-15-15=2.
\]
The corresponding stabilizer subspace has rank \(30\), so the symplectic count gives \(32-30=2\) encoded qudits and Hilbert dimension \(p^2\). The generated
CSV is \path{runs/2026-05-24-symplectic-css-bridge/data/symplectic_css_bridge_summary.csv}.
